% Fermi-LAT analysis of GRB 260924A.
% Numbers come from lat_numbers.json and the tab_lat_*.tex files written by make_lat_products.py.

\subsection{Fermi-LAT data analysis}
\label{sec:lat}

\subsubsection{Data and event selection}
\label{sec:lat:data}

GRB 260924A triggered the \emph{Fermi} Gamma-ray Burst Monitor \citep[GBM;][]{Meegan2009} at
$T_0 = 2026$ September 24 01:13:23.562 UTC (mission elapsed time 811905208.562~s; trigger \texttt{bn260924051}).
We analysed the Large Area Telescope \citep[LAT;][]{Atwood2009} data at the X-ray afterglow position,
$(\alpha, \delta)_{\rm J2000} = (133\fdg76, +12\fdg04)$ \citep{XRTposition}, at the redshift $z = 3.51$ \citep{Redshift}.
We retrieved the extended Pass~8 event file, which contains the transient event classes, and the spacecraft file from
the \emph{Fermi} Science Support Center for $T_0-100$~s to $T_0+1100$~s within $25^\circ$ of the burst position.
The analysis used Fermitools v2.5.1 \citep{Fermitools2019} with the \texttt{P8R3\_TRANSIENT020\_V3} instrument
response functions \citep{Bruel2018}, FRONT and BACK events combined, and an energy range of 100~MeV--100~GeV.
\texttt{TRANSIENT020} is the class recommended for GRBs, where the source dominates the photon count over short
intervals and a larger effective area outweighs a higher background rate.

To suppress gamma rays from the Earth's limb we applied a zenith-angle cut of $\zeta_{\max} = 100^\circ$
together with the region-of-interest (ROI) based time filter of \texttt{gtmktime}. That filter keeps a time $t$
only if the entire circular ROI of radius $r$ lies above the limb,
\begin{equation}
  \zeta(t) + r \le \zeta_{\max},
  \label{eq:roicut}
\end{equation}
where $\zeta(t)$ is the zenith angle of the burst. We also required good data quality and science configuration of
the instrument (\texttt{DATA\_QUAL>0 \&\& LAT\_CONFIG==1}); all spacecraft records in the interval satisfied
these conditions, and the observatory did not enter the South Atlantic Anomaly.

\subsubsection{Definition of the analysis intervals}
\label{sec:lat:windows}

All interval boundaries were derived from the data rather than fixed by hand. They follow from three
constraints: the Earth-limb condition of Eq.~(\ref{eq:roicut}), the off-axis angle of the burst, and the
structure of the prompt keV--MeV emission.

\paragraph{Observing geometry.}
Figure~\ref{fig:lat:geometry} shows the angles of the burst during the observation. At $T_0$ the burst
was $\theta = 57\fdg9$ from the LAT boresight and $\zeta = 88\fdg5$ from the zenith. The zenith angle then
decreased steadily as the burst rose above the limb. By Eq.~(\ref{eq:roicut}), a $10^\circ$ ROI is usable
from $T_0-20.0$~s, whereas a $15^\circ$ ROI becomes usable only from $t_{15} = T_0+70.0$~s. Both stay usable
until the end of the data. The off-axis angle decreased to a minimum of $37\fdg2$ at $T_0+655$~s and then
increased, crossing $\theta_{\max} = 70^\circ$ at $T_0+972.8$~s. Beyond this angle the effective area is
small and the instrument response is less well calibrated, so we ended the LAT window there.

\begin{figure}
  \centering
  \includegraphics[width=\columnwidth]{figures/lat_geometry.pdf}
  \caption{Off-axis angle $\theta$ (solid) and zenith angle $\zeta$ (dashed) of GRB 260924A. Dotted horizontal
  lines mark the largest zenith angle allowed by Eq.~(\ref{eq:roicut}) for the $10^\circ$ (green) and $15^\circ$
  (orange) ROIs, and the dash-dotted line the off-axis limit $\theta_{\max}=70^\circ$. The shaded regions are the
  windows of Analysis~A (green) and Analysis~B (orange).}
  \label{fig:lat:geometry}
\end{figure}

\paragraph{Prompt emission episodes.}
We identified the emission episodes in the GBM CTIME data, which have 0.256~s native resolution, rebinned to
1.024~s. We summed the 10--1000~keV light curves of the three NaI detectors with the largest excess over
background (n3, n5 and n4, at 201, 162 and 160$\sigma$). The background rate was modelled as a polynomial
\begin{equation}
  B(t) = \sum_{k=0}^{n} c_k\, t^k
  \label{eq:gbmbkg}
\end{equation}
fitted to the intervals $[-150, -15]$~s and $[420, 600]$~s. We adopted the lowest order $n$ whose mean squared
Poisson residual in these intervals did not exceed 1.2. That was $n = 3$: the residuals were 1.228, 1.201 and 1.198
for $n = 1$, 2 and 3. We then segmented the light curve with Bayesian blocks \citep{Scargle2013}, using a false
positive probability $p_0 = 0.01$ per change point. For each block $j$ we computed the excess significance
$S_j = (N_j - B_j)/\sqrt{B_j}$ from the observed and background counts, and defined an episode as a maximal run of
consecutive blocks with $S_j \ge 3$ between the two background intervals.

This yielded two episodes: $[-8.2, 231.5]$~s, made of 64 blocks, and $[331.8, 397.3]$~s, made of 15 blocks
(Figure~\ref{fig:lat:lc}a). The first episode was the same for every background order considered. The start of the
second moved to 303.1~s for $n = 1$. For comparison, the GBM catalogue gives $T_{90} = 337.9$~s starting at $T_0+9.2$~s.
After $\sim$400~s, the excess in GBM depends on the assumed background model; across all plausible models it ranges
from $-13\sigma$ to $+15\sigma$. GBM therefore neither detects nor excludes keV emission at those times.

\paragraph{Analyses A and B.}
The geometry offers a trade-off. A $15^\circ$ ROI collects more of the point-spread function (PSF) at
100~MeV, where the 68\% containment radius of \texttt{TRANSIENT020} is several degrees, but it misses the first
70~s of the burst. We therefore performed two independent analyses that differ only in the ROI:
\begin{description}
  \item[Analysis A] used a $10^\circ$ ROI over the first GBM episode, $[-8.2, 231.5]$~s. We fitted the whole episode,
    the part before $t_{15}$ that only this ROI covers, and the interval shared with Analysis~B.
  \item[Analysis B] used a $15^\circ$ ROI over the full LAT window, $[70.0, 972.8]$~s. We fitted it as a whole and
    in adaptive time bins (Section~\ref{sec:lat:bins}).
\end{description}
The common interval $[t_{15}, t_{\rm ep1,end}] = [70.0, 231.5]$~s is covered by both and serves to compare them.

\subsubsection{Likelihood analysis}
\label{sec:lat:like}

For each interval we performed an unbinned maximum-likelihood analysis \citep{Mattox1996} with
\texttt{pyLikelihood}. The log-likelihood is
\begin{equation}
  \ln\mathcal{L} = \sum_{i=1}^{N_{\rm ev}} \ln M(E'_i, \hat{p}'_i, t_i) - N_{\rm pred},
  \label{eq:like}
\end{equation}
with the model count density
\begin{equation}
  M(E', \hat{p}', t) = \sum_s \int\!\!\int R(E', \hat{p}' \mid E, \hat{p}, t)\, S_s(E, \hat{p})\, dE\, d\Omega,
  \label{eq:model}
\end{equation}
where primed quantities are reconstructed and unprimed are true energies and directions. $R$ is the instrument
response, which factorises into the effective area, the PSF and the energy dispersion. $S_s$ is the intensity of
source $s$, and $N_{\rm pred}$ is the total number of counts the model predicts in the ROI. The exposure entering
$N_{\rm pred}$ was computed from a livetime cube with $\Delta\cos\theta = 0.025$ and an exposure map extending
$10^\circ$ beyond the ROI. Energy dispersion was not modelled.

The burst was modelled as a point source with a power-law spectrum parametrised by the integrated photon flux,
\begin{equation}
  \frac{dN}{dE} = F\,\frac{(\gamma + 1)\,E^{\gamma}}{E_{\max}^{\gamma+1} - E_{\min}^{\gamma+1}},
  \qquad
  F = \int_{E_{\min}}^{E_{\max}} \frac{dN}{dE}\,dE,
  \label{eq:pl2}
\end{equation}
with $E_{\min} = 100$~MeV, $E_{\max} = 100$~GeV, and both $F$ and the photon index $\gamma$ free. With this
parametrisation $F$ and $\gamma$ are only weakly correlated. The energy flux is
$F_E = \int E\,(dN/dE)\,dE$. The background model contained the Galactic interstellar emission model
\texttt{gll\_iem\_v07} and the isotropic template \texttt{iso\_P8R3\_TRANSIENT020\_V3\_v1}, each with a free
normalization constrained to $[0, 10]$. It also contained the 243 4FGL-DR4 point sources \citep{Abdollahi2022,
Ballet2023} within $25^\circ$, fixed at their catalogue spectra; no extended sources lie in this region. The fit
used the DRMNFB optimizer followed by \textsc{NewMinuit}. We took parameter uncertainties from the covariance matrix.

The detection significance is the test statistic
\begin{equation}
  {\rm TS} = 2\left(\ln\mathcal{L}_1 - \ln\mathcal{L}_0\right),
  \label{eq:ts}
\end{equation}
where $\mathcal{L}_1$ and $\mathcal{L}_0$ are the maximum likelihoods with and without the burst; the background was
refitted in the null case. By Wilks' theorem \citep{Wilks1938}, TS follows a $\chi^2$ distribution with two degrees
of freedom when no source is present, so ${\rm TS} = 25$ corresponds to $4.6\sigma$. We consider ${\rm TS} \ge 25$
a detection. Otherwise we fixed $\gamma = -2$ and quoted the 95\% one-sided profile-likelihood upper limit on $F$,
where $\ln\mathcal{L}$ falls by $1.355$ from its maximum.

Analyses A and B are not statistically independent. Every event in the $10^\circ$ ROI is also in the $15^\circ$
ROI, so the two fits of the common interval share their photons. We therefore quote the significance of an index
difference $\Delta\gamma = \gamma_{10} - \gamma_{15}$ in two ways: assuming independent errors,
$\sigma_{\rm naive} = (\sigma_{10}^2 + \sigma_{15}^2)^{1/2}$, and for nested samples,
\begin{equation}
  \sigma_{\rm nested} = \left|\sigma_{10}^2 - \sigma_{15}^2\right|^{1/2},
  \label{eq:nested}
\end{equation}
which is the appropriate error when one data set contains the other. The naive form understates the significance
of any difference.

\subsubsection{Adaptive time binning}
\label{sec:lat:bins}

We first cut the Analysis~B window at the GBM episode boundaries inside it (231.5, 331.8 and 397.3~s), so that
no bin mixes episodes. We then bisected each segment at the median arrival time of events within $3^\circ$ of the
burst, keeping the split only if both halves reached ${\rm TS} \ge 25$ and lasted at least 5~s. This was
repeated until no split was accepted. The procedure produced 27 bins, all detections. Table~\ref{tab:lat:bins}
lists them.

To test the robustness of the results to the event class, we refitted the common interval in
\texttt{TRANSIENT020E} and \texttt{SOURCE} for both ROIs. We also refitted every Analysis~B bin in \texttt{SOURCE},
with the \texttt{P8R3\_SOURCE\_V3} responses and the matching isotropic template. Finally, we computed with
\texttt{gtsrcprob} the probability that each event belongs to the burst,
\begin{equation}
  p_i = \frac{M_{\rm GRB}(E'_i, \hat{p}'_i)}{\sum_s M_s(E'_i, \hat{p}'_i)},
  \label{eq:prob}
\end{equation}
using the best-fit models of the full Analysis~A and B windows.

\subsubsection{Results}
\label{sec:lat:results}

\begin{table*}
  \centering
  \caption{Likelihood results for the analysis windows. $N_{\rm ev}$ is the number of events in the ROI and
  $N_{\rm pred}$ the number predicted for the burst. $F$ and $F_E$ are the 0.1--100~GeV photon and energy fluxes.
  The interval before $t_{15}$ is not detected; we give the 95\% upper limit with $\gamma=-2$ fixed.}
  \label{tab:lat:windows}
  \begin{tabular}{llrrrrccc}
    \hline
    Analysis & ROI & Interval & $N_{\rm ev}$ & $N_{\rm pred}$ & TS & $\gamma$ & $F$ & $F_E$ \\
     & & (s) & & & & & ($10^{-5}$ ph cm$^{-2}$ s$^{-1}$) & ($10^{-8}$ erg cm$^{-2}$ s$^{-1}$) \\
    \hline
    A & $10^\circ$ & $-8.2$--$231.5$ & 228 & 228 & 1316 & $-2.22\pm0.07$ & $29.21\pm1.98$ & $20.39\pm2.79$ \\
A & $10^\circ$ & $-8.2$--$70.0$ & 6 & 5 & 18 & $-2$ (fixed) & $<4.29$ & \ldots \\
A & $10^\circ$ & $70.0$--$231.5$ & 222 & 222 & 1411 & $-2.21\pm0.07$ & $40.05\pm2.74$ & $28.33\pm3.97$ \\
B & $15^\circ$ & $70.0$--$231.5$ & 237 & 229 & 1424 & $-2.21\pm0.08$ & $39.83\pm2.78$ & $28.14\pm3.96$ \\
B & $15^\circ$ & $70.0$--$972.8$ & 385 & 346 & 1515 & $-2.21\pm0.06$ & $9.17\pm0.54$ & $6.49\pm0.76$ \\
\hline

  \end{tabular}
\end{table*}

\paragraph{Detection and onset.}
The burst is strongly detected in both analyses (Table~\ref{tab:lat:windows}). Over the first GBM episode the
$10^\circ$ analysis gives ${\rm TS} = 1316$, $\gamma = -2.22 \pm 0.07$ and
$F = (2.92 \pm 0.20)\times10^{-4}$~ph~cm$^{-2}$~s$^{-1}$. Over the full LAT window the $15^\circ$ analysis gives
${\rm TS} = 1515$, $\gamma = -2.21 \pm 0.06$ and $F = (9.17 \pm 0.54)\times10^{-5}$~ph~cm$^{-2}$~s$^{-1}$.
Between $-8.2$ and $70.0$~s, the stretch only the $10^\circ$ ROI covers, we find just six events with $p > 0.9$
(Figure~\ref{fig:lat:lc}b) and ${\rm TS} = 18$. The 95\% upper limit is
$F < 4.3\times10^{-5}$~ph~cm$^{-2}$~s$^{-1}$, a factor of 36 below the flux in the first $15^\circ$ bin,
$70.0$--$82.8$~s, where $F = (1.53 \pm 0.20)\times10^{-3}$~ph~cm$^{-2}$~s$^{-1}$. The upper limit is corrected
for exposure. If emission at the flux of the first $15^\circ$ bin had occurred before $t_{15}$, it could have
lasted at most $\sim$2~s. The bright $>100$~MeV emission therefore started close to $t_{15}$, at least
$\sim$78~s (17~s in the rest frame) after the GBM episode began.

\begin{figure*}
  \centering
  \includegraphics[width=0.85\textwidth]{figures/lat_lightcurve.pdf}
  \caption{Multi-wavelength light curve of GRB 260924A. (a) Background-subtracted GBM NaI count rate (grey) and its
  Bayesian-block representation (black). (b) Energies of LAT events in the $15^\circ$ ROI, with burst probability
  $p>0.9$ (black) or $p<0.9$ (grey), and $p>0.9$ events in the $10^\circ$ ROI before $t_{15}$ (green squares).
  (c) 0.1--100~GeV photon flux in the Analysis~B bins (orange) and Analysis~A windows (green); the arrow is the 95\%
  upper limit before $t_{15}$. (d) Photon index; the orange band is the weighted mean of the Analysis~B bins with its
  $1\sigma$ error. Grey shading marks the GBM episodes; vertical dotted lines mark the start and end of the
  Analysis~B window.}
  \label{fig:lat:lc}
\end{figure*}

\paragraph{Comparison of the two ROIs.}
In the common interval, the two ROIs give practically identical spectra (Table~\ref{tab:lat:common};
Figure~\ref{fig:lat:sed}a). With \texttt{TRANSIENT020} we obtain $\gamma = -2.21 \pm 0.07$ in $10^\circ$ and
$-2.21 \pm 0.08$ in $15^\circ$. The difference is $\Delta\gamma = 0.001$, which is $0.13\sigma$ even with the
nested error of Eq.~(\ref{eq:nested}), and the fluxes agree to 0.5\%. The 15 extra events in the larger ROI are
absorbed almost entirely by the diffuse components: $N_{\rm pred}$ for the burst rises from 222 to 229. The same
holds for \texttt{TRANSIENT020E} and for \texttt{SOURCE} (Figure~\ref{fig:lat:sed}b). \texttt{SOURCE} gives a
slightly harder index, $\gamma \approx -2.15$, which is consistent with \texttt{TRANSIENT020} within $0.6\sigma$.

\begin{table*}
  \centering
  \caption{Spectral comparison of the $10^\circ$ (Analysis~A) and $15^\circ$ (Analysis~B) ROIs on the common interval
  $[70.0, 231.5]$~s for three event classes. Entries with a slash give the $10^\circ$ / $15^\circ$ values. The
  significance of $\Delta\gamma=\gamma_{10}-\gamma_{15}$ is quoted with independent and nested
  (Eq.~\ref{eq:nested}) errors.}
  \label{tab:lat:common}
  \begin{tabular}{lcccccc}
    \hline
    Event class & $N_{\rm ev}$ & TS & $\gamma_{10}$ & $\gamma_{15}$ & $\Delta\gamma/\sigma$ (naive / nested) & $F_{10}/F_{15}$ \\
    \hline
    \texttt{TRANSIENT020E} & 236 / 251 & 1510 / 1523 & $-2.21\pm0.07$ & $-2.21\pm0.08$ & $+0.01$ / $+0.03$ & 1.003 \\
\texttt{TRANSIENT020} & 222 / 237 & 1411 / 1424 & $-2.21\pm0.07$ & $-2.21\pm0.08$ & $+0.01$ / $+0.13$ & 1.005 \\
\texttt{SOURCE} & 144 / 146 & 1128 / 1139 & $-2.16\pm0.08$ & $-2.15\pm0.08$ & $-0.03$ / $-0.36$ & 1.009 \\
\hline

  \end{tabular}
\end{table*}

\begin{figure*}
  \centering
  \includegraphics[width=0.85\textwidth]{figures/lat_sed_common.pdf}
  \caption{(a) Spectral energy distribution of the common interval $[70.0, 231.5]$~s in \texttt{TRANSIENT020} for the
  $10^\circ$ and $15^\circ$ ROIs. The bands are the 68\% range of $10^4$ power laws drawn from the fitted $(F,\gamma)$
  covariance. (b) Photon index on the common interval for the three event classes and both ROIs.}
  \label{fig:lat:sed}
\end{figure*}

\paragraph{Spectral evolution.}
The photon index shows no significant change across the 27 bins (Figure~\ref{fig:lat:lc}d). A constant index
fits with weighted mean $\bar{\gamma} = -2.18 \pm 0.06$ and $\chi^2 = 21.2$ for 26 degrees of freedom ($p = 0.73$).
The 14 bins inside the first GBM episode average $-2.18 \pm 0.07$ and the 13 later bins $-2.19 \pm 0.11$. The
largest deviation is the bin $141.6$--$151.9$~s, with $\gamma = -1.63 \pm 0.21$, $2.6\sigma$ harder than the
mean. With 27 bins, a deviation at least that large occurs by chance with probability 0.22. The \texttt{SOURCE}
refits confirm this picture. Of the 27 bins, 20 remain detected, and their indices agree with \texttt{TRANSIENT020}
to better than $0.8\sigma$ in every case, with a median flux ratio of 0.94. The remaining seven have
${\rm TS} = 20$--25 in \texttt{SOURCE}, as expected for a class that keeps about two thirds of the photons.

\paragraph{Spectral curvature.}
To test whether the spectrum is curved rather than a single power law, we refitted the main intervals with a
log-parabola,
\begin{equation}
  \frac{dN}{dE} = N_0 \left(\frac{E}{E_b}\right)^{-\alpha - \beta_{\rm c}\ln(E/E_b)},
  \label{eq:logparabola}
\end{equation}
with the pivot fixed at $E_b = 1$~GeV and $N_0$, $\alpha$ and $\beta_{\rm c}$ free. We write the curvature as
$\beta_{\rm c}$ because $\beta$ is used below for the energy spectral index. Setting $\beta_{\rm c} = 0$ recovers a
power law, so the improvement
\begin{equation}
  {\rm TS}_{\rm curv} = 2\left(\ln\mathcal{L}_{\rm LP} - \ln\mathcal{L}_{\rm PL}\right)
  \label{eq:tscurv}
\end{equation}
has one degree of freedom, and $\sqrt{{\rm TS}_{\rm curv}}$ is its significance. A positive $\beta_{\rm c}$ means the
spectrum softens toward higher energy, as a cooling break or a cutoff inside the band would, and a negative value
means it hardens. The results are in Table~\ref{tab:lat:curv}. In every interval $\beta_{\rm c}$ is slightly
negative but consistent with zero: the largest improvement is ${\rm TS}_{\rm curv} = 0.99$ ($1.0\sigma$) for the
full Analysis~B window, where $\beta_{\rm c} = -0.037 \pm 0.036$. The interval after the first GBM episode, the part
of the emission most likely to be afterglow, gives $\beta_{\rm c} = -0.049 \pm 0.063$ and
${\rm TS}_{\rm curv} = 0.53$. Because $\beta_{\rm c}$ is not positive, the spectral energy distribution has no peak
inside 0.1--100~GeV. A power law is therefore an adequate description of the LAT spectrum, and there is no evidence
for a spectral break within the band.

\begin{table*}
  \centering
  \caption{Power-law and log-parabola fits (Eq.~\ref{eq:logparabola}) of the main intervals, all in
  \texttt{TRANSIENT020}. $\gamma$ is the power-law photon index; $\alpha$ and $\beta_{\rm c}$ are the log-parabola
  parameters at $E_b=1$~GeV. ${\rm TS}_{\rm curv}$ (Eq.~\ref{eq:tscurv}) is the improvement of the log-parabola over
  the power law.}
  \label{tab:lat:curv}
  \begin{tabular}{llccccc}
    \hline
    Analysis & ROI & Interval (s) & $\gamma$ & $\alpha$ & $\beta_{\rm c}$ & ${\rm TS}_{\rm curv}$ \\
    \hline
    A & $10^\circ$ & $-8.2$--$231.5$ & $-2.22\pm0.07$ & $2.18\pm0.08$ & $-0.033\pm0.043$ & 0.52 \\
A & $10^\circ$ & $70.0$--$231.5$ & $-2.21\pm0.07$ & $2.17\pm0.08$ & $-0.036\pm0.043$ & 0.63 \\
B & $15^\circ$ & $70.0$--$231.5$ & $-2.21\pm0.08$ & $2.17\pm0.08$ & $-0.039\pm0.043$ & 0.76 \\
B & $15^\circ$ & $70.0$--$972.8$ & $-2.21\pm0.06$ & $2.17\pm0.07$ & $-0.037\pm0.036$ & 0.99 \\
B & $15^\circ$ & $231.5$--$972.8$ & $-2.27\pm0.11$ & $2.21\pm0.12$ & $-0.049\pm0.063$ & 0.53 \\
\hline

  \end{tabular}
\end{table*}

\begin{table*}
  \centering
  \caption{Analysis~B: $15^\circ$ ROI, \texttt{TRANSIENT020}, adaptive time bins. Columns 2--7 give the bin
  interval, number of events, TS, photon index, and 0.1--100~GeV photon and energy fluxes (units as in
  Table~\ref{tab:lat:windows}). The last two columns give the index and TS of the \texttt{SOURCE}-class refit;
  ``\ldots'' marks bins with ${\rm TS}<25$ in that class.}
  \label{tab:lat:bins}
  \begin{tabular}{rrrrcccrr}
    \hline
    Bin & Interval (s) & $N_{\rm ev}$ & TS & $\gamma$ & $F$ & $F_E$ & $\gamma_{\rm SOURCE}$ & TS$_{\rm SOURCE}$ \\
    \hline
    1 & $70.0$--$82.8$ & 61 & 490 & $-2.43\pm0.17$ & $153.26\pm19.96$ & $77.03\pm17.31$ & $-2.32\pm0.20$ & 325 \\
2 & $82.8$--$91.6$ & 29 & 226 & $-2.22\pm0.21$ & $102.67\pm19.39$ & $70.97\pm27.55$ & $-2.21\pm0.22$ & 212 \\
3 & $91.6$--$98.6$ & 17 & 94 & $-2.39\pm0.33$ & $70.63\pm18.92$ & $37.60\pm17.58$ & $-2.57\pm0.38$ & 83 \\
4 & $98.6$--$105.5$ & 15 & 103 & $-2.14\pm0.27$ & $61.95\pm16.26$ & $50.58\pm29.15$ & $-2.02\pm0.29$ & 90 \\
5 & $105.5$--$110.6$ & 12 & 80 & $-2.39\pm0.36$ & $67.60\pm19.85$ & $36.17\pm19.13$ & $-2.65\pm0.45$ & 79 \\
6 & $110.6$--$125.1$ & 18 & 104 & $-2.02\pm0.23$ & $34.02\pm8.15$ & $35.61\pm19.93$ & $-1.75\pm0.29$ & 64 \\
7 & $125.1$--$141.6$ & 24 & 170 & $-2.30\pm0.25$ & $41.43\pm8.59$ & $25.11\pm10.13$ & $-2.20\pm0.27$ & 140 \\
8 & $141.6$--$151.8$ & 11 & 112 & $-1.63\pm0.21$ & $27.04\pm8.27$ & $89.75\pm64.02$ & $-1.53\pm0.23$ & 103 \\
9 & $151.8$--$163.8$ & 8 & 32 & $-1.94\pm0.38$ & $14.65\pm6.35$ & $19.05\pm18.58$ & \ldots & 22 \\
10 & $163.8$--$171.6$ & 6 & 28 & $-2.21\pm0.52$ & $17.15\pm8.38$ & $12.23\pm12.09$ & \ldots & 20 \\
11 & $171.6$--$194.1$ & 12 & 35 & $-2.32\pm0.44$ & $11.45\pm4.83$ & $6.71\pm4.34$ & \ldots & 23 \\
12 & $194.1$--$200.4$ & 6 & 42 & $-2.08\pm0.42$ & $24.01\pm9.94$ & $22.08\pm20.93$ & $-2.28\pm0.44$ & 50 \\
13 & $200.4$--$209.7$ & 7 & 33 & $-2.02\pm0.37$ & $18.77\pm7.20$ & $19.84\pm17.92$ & $-1.91\pm0.37$ & 27 \\
14 & $209.7$--$231.5$ & 11 & 51 & $-2.59\pm0.44$ & $13.79\pm4.21$ & $5.86\pm2.77$ & \ldots & 23 \\
15 & $231.5$--$258.9$ & 11 & 39 & $-2.47\pm0.41$ & $10.44\pm3.20$ & $5.05\pm2.65$ & $-2.61\pm0.55$ & 28 \\
16 & $258.9$--$281.6$ & 6 & 27 & $-2.23\pm0.51$ & $5.33\pm2.48$ & $3.67\pm3.49$ & \ldots & 25 \\
17 & $281.6$--$331.8$ & 14 & 62 & $-2.10\pm0.28$ & $6.64\pm1.80$ & $5.78\pm3.55$ & $-2.26\pm0.34$ & 50 \\
18 & $331.8$--$356.4$ & 9 & 51 & $-1.88\pm0.36$ & $6.31\pm2.55$ & $9.68\pm10.11$ & $-1.68\pm0.40$ & 33 \\
19 & $356.4$--$397.3$ & 11 & 33 & $-3.01\pm0.68$ & $4.57\pm1.79$ & $1.45\pm0.67$ & $-2.94\pm0.66$ & 42 \\
20 & $397.3$--$441.2$ & 9 & 38 & $-2.69\pm0.58$ & $3.82\pm1.48$ & $1.49\pm0.84$ & $-2.39\pm0.58$ & 26 \\
21 & $441.2$--$480.9$ & 9 & 49 & $-1.67\pm0.28$ & $3.35\pm1.31$ & $9.64\pm8.95$ & $-1.42\pm0.30$ & 40 \\
22 & $480.9$--$563.4$ & 17 & 38 & $-2.30\pm0.41$ & $2.55\pm0.99$ & $1.55\pm0.99$ & $-2.80\pm0.47$ & 41 \\
23 & $563.4$--$609.9$ & 12 & 54 & $-2.24\pm0.36$ & $4.73\pm1.59$ & $3.18\pm2.02$ & $-2.02\pm0.35$ & 53 \\
24 & $609.9$--$669.5$ & 13 & 38 & $-2.10\pm0.39$ & $2.54\pm1.09$ & $2.24\pm1.88$ & $-2.52\pm0.43$ & 46 \\
25 & $669.5$--$692.9$ & 9 & 38 & $-2.59\pm0.49$ & $7.55\pm2.54$ & $3.22\pm1.73$ & \ldots & 24 \\
26 & $692.9$--$763.2$ & 9 & 48 & $-2.32\pm0.41$ & $2.58\pm0.88$ & $1.52\pm1.01$ & $-2.19\pm0.46$ & 40 \\
27 & $763.2$--$972.8$ & 19 & 31 & $-2.67\pm0.47$ & $1.91\pm0.68$ & $0.75\pm0.33$ & \ldots & 22 \\
\hline

  \end{tabular}
\end{table*}

\paragraph{Temporal decay.}
The LAT flux peaks in the first available bin and then declines over an order of magnitude in time
(Figure~\ref{fig:lat:decay}). We fitted the bin fluxes, evaluated at the logarithmic bin centre
$t_c = (t_{\rm start}t_{\rm stop})^{1/2}$, with a single power law $F(t) \propto t^{-\alpha}$ and with a continuous
broken power law,
\begin{equation}
  F(t) = F_b \times
  \begin{cases}
    (t/t_b)^{-\alpha_1}, & t < t_b,\\
    (t/t_b)^{-\alpha_2}, & t \ge t_b.
  \end{cases}
  \label{eq:bpl}
\end{equation}
The single power law gives $\alpha = 1.80 \pm 0.07$ with $\chi^2 = 42.3$ for 25 degrees of freedom, a poor fit
($p = 0.017$). The broken power law gives $\alpha_1 = 2.72 \pm 0.30$, $\alpha_2 = 1.22 \pm 0.18$ and
$t_b = 175^{+31}_{-27}$~s, with $\chi^2 = 22.2$ for 23 degrees of freedom. The improvement,
$\Delta\chi^2 = 20.1$ for two extra parameters, is significant: an $F$-test gives $p = 6\times10^{-4}$.
Fitting the segments separately gives the same result. Inside the first GBM episode the flux decays as
$t^{-2.37 \pm 0.20}$, and after it ends as $t^{-0.92 \pm 0.25}$. The break lies within the first GBM episode,
$\sim$55~s before its end. Restricting the fit to bins after the second episode leaves
$\alpha = 0.46 \pm 0.58$, which is not constraining.

\begin{figure}
  \centering
  \includegraphics[width=\columnwidth]{figures/lat_decay.pdf}
  \caption{0.1--100~GeV photon flux of the Analysis~B bins at their logarithmic centres, with the best-fitting single
  (dashed) and broken (solid) power laws of Eq.~(\ref{eq:bpl}); the dotted line marks $t_b$. Grey bands are the GBM
  episodes.}
  \label{fig:lat:decay}
\end{figure}

\paragraph{Energetics.}
We integrated the energy flux over each window to obtain the observed fluence $S = F_E\,\Delta t$, and converted it
to an isotropic-equivalent energy in the rest-frame band 0.1--100~GeV,
\begin{equation}
  E_{\rm iso} = \frac{4\pi d_L^2}{1+z}\, S\, k,
  \qquad
  k = \frac{\int_{E_1/(1+z)}^{E_2/(1+z)} E^{\gamma+1}\, dE}{\int_{E_{\min}}^{E_{\max}} E^{\gamma+1}\, dE},
  \label{eq:eiso}
\end{equation}
where $E_1 = 100$~MeV and $E_2 = 100$~GeV \citep{Bloom2001}. We used $d_L = 9.69\times10^{28}$~cm, from the
\citet{Planck2020} cosmology. For $\gamma \approx -2.2$ the correction is $k = 1.39 \pm 0.16$, where the error comes
from the index uncertainty. The correction requires extrapolating the observed power law down to 22~MeV. Over the first
GBM episode, $S = (4.88 \pm 0.67)\times10^{-5}$~erg~cm$^{-2}$ and
$E_{\rm iso} = (1.8 \pm 0.3)\times10^{54}$~erg. Over the full LAT window, $S = (5.86 \pm 0.69)\times10^{-5}$~erg~cm$^{-2}$
and $E_{\rm iso} = (2.1 \pm 0.3)\times10^{54}$~erg. The two windows overlap, so these are not additive. The
corresponding mean isotropic luminosities in the rest frame are $3.3\times10^{52}$ and $1.0\times10^{52}$~erg~s$^{-1}$.

\paragraph{High-energy photons.}
In the Analysis~B model, 295 events have $p > 0.9$; of these, 28 are above 1~GeV and four above 10~GeV
(Table~\ref{tab:lat:photons}). The highest-energy photon, 43.3~GeV, arrived at $T_0+82.24$~s, $0\fdg07$ from the
burst position and at $\theta = 54\fdg9$, with $p > 0.9999$ in both analyses. It corresponds to 195~GeV in the
rest frame. Further photons of 24.9~GeV (at 122.9~s) and 11.6~GeV (at 141.6~s) fall within the first GBM episode.
A 14.2~GeV photon arrived at 304.1~s, during the quiescent interval between the two GBM episodes. The photon of
5.5~GeV at 906.7~s lies $5\fdg7$ from the burst and has $p = 0.05$, so it is attributed to the background.

\begin{table}
  \centering
  \caption{Highest-energy LAT photons associated with GRB 260924A ($p>0.9$). $E_{\rm rest}=E(1+z)$; $\rho$ is the
  angular distance from the burst position; $p_A$ and $p_B$ are the probabilities in the Analysis~A and B models
  (``\ldots'': outside the Analysis~A window).}
  \label{tab:lat:photons}
  \begin{tabular}{rrrrrcc}
    \hline
    $t-T_0$ & $E$ & $E_{\rm rest}$ & $\rho$ & $\theta$ & $p_A$ & $p_B$ \\
    (s) & (GeV) & (GeV) & (deg) & (deg) & & \\
    \hline
    82.24 & 43.3 & 195 & 0.07 & 54.9 & 1.000 & 1.000 \\
122.92 & 24.9 & 113 & 0.41 & 53.8 & 1.000 & 0.999 \\
304.15 & 14.2 & 64 & 0.08 & 47.9 & \ldots & 1.000 \\
141.64 & 11.6 & 52 & 0.19 & 53.1 & 1.000 & 1.000 \\
665.55 & 5.3 & 24 & 0.06 & 37.7 & \ldots & 1.000 \\
98.99 & 4.0 & 18 & 0.29 & 54.3 & 1.000 & 1.000 \\
88.12 & 3.5 & 16 & 0.60 & 54.3 & 1.000 & 0.999 \\
587.17 & 3.1 & 14 & 0.12 & 37.9 & \ldots & 1.000 \\
136.79 & 2.6 & 12 & 0.24 & 52.9 & 1.000 & 1.000 \\
147.86 & 2.4 & 11 & 0.23 & 52.6 & 1.000 & 1.000 \\
\hline

  \end{tabular}
\end{table}

\subsubsection{Interpretation}
\label{sec:lat:discussion}

The LAT results are robust to the analysis choices we tested. Enlarging the ROI from $10^\circ$ to $15^\circ$
changes neither the spectrum nor the flux in the common interval, so the source dominates the photon count and the
background model has little effect. The extra $5^\circ$ annulus mostly adds background, which the fit attributes
correctly to the diffuse components. The agreement among \texttt{TRANSIENT020E}, \texttt{TRANSIENT020} and
\texttt{SOURCE} likewise excludes a significant contribution from residual cosmic rays, to which the transient
classes are most exposed. The only gain of the $10^\circ$ analysis is temporal coverage of $-8$ to 70~s, which shows
that the GeV emission was faint before $\sim$70~s.

GRB 260924A shows the behaviour common to bright LAT bursts \citep{Ajello2019}. The $>100$~MeV emission rises after
the keV--MeV emission has started, peaks during the bright prompt phase, and lasts well beyond it. The LAT
detected the burst in every bin up to the $\theta_{\max}$ cut at 973~s, about 2.8 times $T_{90}$, while GBM has
no reliable detection after $\sim$400~s. The isotropic-equivalent energy of $\sim2\times10^{54}$~erg in the
0.1--100~GeV band places it at the high end of the distribution for LAT-detected bursts \citep{Ajello2019}.

Two observations together point to a change of emission regime near the end of the first GBM episode. The photon
index remains $\bar{\gamma} = -2.18 \pm 0.06$ throughout, but the decay steepens from
$\alpha_1 \approx 2.7$ to $\alpha_2 \approx 1.2$ at $t_b \approx 175$~s. A decay as steep as
$\alpha_1 \approx 2.7$, while the keV emission is still bright, suggests that a component tied to the prompt
phase contributes to the early LAT flux. This could be internal emission, or an external shock still being
energised, with the time origin placed at $T_0$ rather than at the onset of the afterglow. The late decay can be
compared with synchrotron emission from an adiabatic forward shock \citep[e.g.][]{Kumar2009, Kumar2010}.
With spectral index $\beta = -\bar{\gamma} - 1 = 1.18 \pm 0.06$ and the LAT band above the cooling frequency,
$\beta = p/2$ gives an electron index $p = 2.36 \pm 0.12$. The closure relation then predicts
\begin{equation}
  \alpha = \frac{3p - 2}{4} = \frac{3\beta}{2} - \frac{1}{2} = 1.27 \pm 0.09,
  \label{eq:closure}
\end{equation}
in agreement with $\alpha_2 = 1.22 \pm 0.18$. Placing the band between the injection and cooling frequencies instead
gives $p = 3.36$ and $\alpha = 3\beta/2 = 1.77 \pm 0.09$, $2.7\sigma$ away from $\alpha_2$. The late LAT emission is
therefore consistent with an adiabatic forward shock radiating above the cooling frequency, with $p \approx 2.4$.
The steady photon index, and the absence of curvature in the log-parabola fits
(Table~\ref{tab:lat:curv}), are consistent with the cooling frequency staying below 100~MeV throughout. A combined fit with the
X-ray and optical afterglow (Section~\ref{sec:afterglow}) is needed to test this further.

\paragraph{Caveats.}
Several limitations affect these results. First, the burst was observed $50^\circ$--$58^\circ$ off-axis
during its brightest phase. At those angles the effective area and energy reconstruction are less certain than
on-axis. This mainly affects the energies of individual photons, such as the 43.3~GeV event, not the fitted indices.
Second, we did not model energy dispersion, which can bias power-law indices slightly for spectra dominated by
photons below 300~MeV. Third, in many short bins one of the diffuse normalizations reached a limit of its allowed
range, which is expected when the background contributes only a few photons. Fourth, \textsc{NewMinuit} returned a
non-zero status code for the full Analysis~B window and for the $15^\circ$ \texttt{TRANSIENT020E} fit, so their
uncertainties should be treated with caution; their best-fit values agree with the neighbouring fits. Fifth, the GBM episode
boundaries depend weakly on the background model, as described in Section~\ref{sec:lat:windows}; this affects only
where the Analysis~B segments start and end. Finally, the decay indices are measured with respect to $T_0$. A
different time origin would change $\alpha_1$ much more than $\alpha_2$. We have not derived a bulk Lorentz-factor
limit from the 43.3~GeV photon, because that requires the target photon spectrum and variability time of the
prompt emission.
